\chapter{Vitamin B1 (Thiamine)}

\section*{What Is This Test?}
Vitamin B1, or thiamine, testing assesses thiamine status. Thiamine is converted to thiamine pyrophosphate (TPP), a cofactor required for carbohydrate metabolism, energy production, and normal nerve and cardiac function. Severe deficiency can cause beriberi or Wernicke encephalopathy.

\section*{Alternative Names}
Thiamine level; vitamin B1; thiamine pyrophosphate; TPP; erythrocyte transketolase activity

\section*{Principle}
Testing may measure whole-blood thiamine or thiamine diphosphate, usually by high-performance liquid chromatography or liquid chromatography--mass spectrometry. Erythrocyte transketolase activity, with or without added TPP, is an older functional method and is less commonly used.

\section*{Why This Test Is Ordered}
\begin{itemize}
  \item Suspected Wernicke encephalopathy or beriberi
  \item Unexplained confusion, ataxia, ophthalmoplegia, neuropathy, muscle weakness, or high-output heart failure in a patient at risk of deficiency
  \item Alcohol-use disorder, severe malnutrition, prolonged vomiting, malabsorption, or bariatric surgery
  \item Monitoring selected patients receiving long-term nutrition support or replacement therapy
\end{itemize}

\section*{Primary vs Secondary Test}
Thiamine testing is a supportive test. Suspected Wernicke encephalopathy is a clinical emergency and treatment should not be delayed while waiting for the result.

\section*{How This Test Is Performed}
Whole blood is collected in a laboratory-specified trace-element or EDTA tube. The sample is protected from light and transported promptly because thiamine is unstable.

\section*{What Preparation Is Needed}
Fasting is usually unnecessary. Document recent vitamin supplements, alcohol intake, nutritional support, and thiamine treatment. The laboratory's collection requirements should be followed exactly.

\section*{Risks}
Venipuncture only.

\section*{Understanding Results}
\begin{itemize}
  \item \textbf{Low thiamine:} Supports deficiency but should be interpreted with the clinical picture; a normal result does not reliably exclude tissue deficiency in an acutely ill patient.
  \item \textbf{Risk factors:} Alcohol-use disorder, inadequate food intake, persistent vomiting, malabsorption, dialysis, and prolonged carbohydrate-rich intravenous nutrition without adequate vitamin replacement.
  \item \textbf{Treatment response:} Clinical treatment is guided by symptoms and risk, not by waiting for laboratory confirmation. Neurologic injury may become permanent if treatment is delayed.
\end{itemize}

\section*{Follow-up Tests}
\begin{itemize}
  \item Complete blood count, electrolytes, glucose, magnesium, phosphate, and liver function tests
  \item Assessment for other nutritional deficiencies, including folate, vitamin B12, and vitamin C when malnutrition is present
  \item Neurologic or cardiac assessment when Wernicke encephalopathy or beriberi is suspected
\end{itemize}

\clearpage
\section*{Regulatory Status and Authorization Date}
Regulatory status applies to the specific assay, instrument, manufacturer, and intended use rather than to the test name alone. No single approval date can be assigned to this test category. For a commercial assay, verify the current product in the relevant regulator's database: the U.S. Food and Drug Administration (FDA) clearance, approval, or authorization; India's Central Drugs Standard Control Organisation (CDSCO) licence or permission; the European Union In Vitro Diagnostic Regulation (EU IVDR) CE certificate issued through the applicable conformity-assessment route; or the United Kingdom Medicines and Healthcare products Regulatory Agency (MHRA) registration and UKCA/CE status. Laboratory-developed tests may not have an individual FDA, CDSCO, EU IVDR, or MHRA approval date.
\clearpage
